\section{计算核验、适用范围与后续问题}\label{sec:computation}

\subsection{一般证明与有限检查的分工}
三行主定理的第一条依赖链为：列规则的完备性、首次回返块分类、
唯一分解与对称字作用、固定字计数、轨道平均。
第二条依赖链为：局部部分铺法模型、半棋盘补全双射、两维递推、轨道平均。
两条路线的组合论证互不依赖；它们共享局部覆盖规则、轨道平均和最终代数化简。
方法之间的交叉核对不等同于独立作者的审阅。

原基本块程序对 $n=0,\ldots,8$ 的全部三行铺法作分解、重建及四种对称作用比较，
共检查 $40545$ 个对象；它复用了此前的行优先格点枚举器。
后续半棋盘程序使用新写的列优先顶点编号及最少可配邻居选择规则，
不导入旧枚举器或基本块代码。其完整枚举范围见表 \ref{tab:checks}。

\begin{table}[htbp]
\centering
\caption{半棋盘程序实际完成的全部铺法枚举范围。}\label{tab:checks}
\begin{tabular}{ccl}
\toprule
高度 & 半长度 $n$ & 检查重点\\
\midrule
1 & $0,\ldots,8$ & 最小高度、非忠实群作用\\
3 & $0,\ldots,8$ & 原公式、精确稳定子、与旧报告对照\\
5 & $0,\ldots,4$ & 五行显式公式与中线补全\\
7 & $0,\ldots,2$ & 更高奇数行的双射实例\\
\bottomrule
\end{tabular}
\end{table}

半棋盘枚举器在这些范围合计遍历 $57468$ 个铺法。
每个铺法都计算四种格点对称像和轨道规范代表；
对左右反射或半周旋转固定者，限制到左半、记录末端半块集合并重建全棋盘，
检查重建一致性和按每个末端状态分类的计数。
三行精确稳定子数也与表 \ref{tab:stabilizers} 一致。

符号计算从完整三行、五行转移矩阵求解有理函数，
核对矩阵逆行向量满足对应线性方程且残差为零。
五行六维模型另外核对式 \eqref{eq:certificate} 以及它与完整矩阵模型的一致性。
三、五行有理式系数与整数矩阵迭代在 $n=0,\ldots,100$ 一致；
这一项不是穷举该范围的全部铺法。
上述为 2026 年 10 月 25 日既有研究阶段实际运行的检查，
使用 Python 3.13.0、SymPy 1.13.1 和 Windows，
全部为确定性精确计算，无随机种子或浮点拟合。
原始输出、源码哈希与命令仍保存在附录列出的研究材料中。

有限枚举检查实现与小规模实例；精确多项式恒等式检查矩阵求值；
一般覆盖性与对称解释由本文双射证明承担。
这些证据各有范围，不将有限前缀外推为任意 $n$ 的证明。

\subsection{结论的范围和后续工作}
本文完整推导了三行的目标有理式及递推，
给出精确稳定子分类、渐近比例和任意固定奇数高度的统一矩阵公式。
五行有理式由显式有限代数证书确定。
下一项工作是由独立读者逐步审阅关键双射，尤其是跨中线砖的补全和五行状态压缩，
并继续追踪对称铺法计数的既有文献以确定成果定位。
本文不因 OEIS 页面保留猜想标签而主张全球开放状态，
也不因本地两条推导一致而主张已获得外部审稿确认。

偶数高度和正方形是自然的进一步方向，但分别涉及不同的上下反射边界条件和更大的对称群。
本文没有对这些范围给出结论。精确稳定子、一般奇数高度、五行代数证书、
渐近分析和基本块证明的后续本地形式化属于独立阶段，应另行报告和审阅。
本版的语义审阅包仅覆盖第 \ref{sec:formalization} 节说明的 M1--M6 三行主计数范围。
